% GENERATED FILE. Edit canonical lessons, then run npm run generate:books.
% canonical-source-hash: fnv1a64:341bff92cd875210
% canonical-lessons: MR-C431-bighad, MR-C431-nikami, MR-C431-tutne, MR-C431-phutne, MR-C431-phatne

\chapter{When Things Break Down}
\label{ch:mr-when-things-break-down}
\hlchaptermodality{431}

\begin{chapteropening}
\textbf{By the end of this chapter:} \emph{I can talk about a fault, things out of action, and say something snapped, burst or tore.}

Complete the last lesson of chapter 431: I can talk about a fault, things out of action, and say something snapped, burst or tore.
\end{chapteropening}

\section[\texorpdfstring{\mr{बिघाड} (bighāḍ)}{बिघाड (bighāḍ)}]{\mr{बिघाड} (bighāḍ) — a breakdown; a fault}

\label{lesson:MR-C431-bighad}



\begin{quote}
Before the new one: say the Marathi for a share, then the Marathi for to be left over.
\end{quote}

\subsection*{You'll want to know: \mr{बिघाड}}
\textbf{\mr{बिघाड}} — \emph{bighāḍ} — \textquotedblleft{}a breakdown; a fault\textquotedblright{}.

\begin{quote}
\mr{इंजिनात} \mr{बिघाड} \mr{झाल्यामुळे} \mr{गाडी} \mr{थांबली}.
\end{quote}

\begin{quote}
\emph{iñjināt bighāḍ jhālyāmuḷe gāḍī thāmblī.}
\end{quote}

\begin{quote}
The train stopped because of a fault in the engine.
\end{quote}

Masculine.

\mr{बिघाड} \mr{होणे}, a fault occurs; the place of the fault takes -\mr{त}. The verb is \mr{बिघडणे}, to go wrong.

\begin{cousinweb}
\textbf{\mr{बिघाड}} (\emph{bighāḍ}) is inherited, from the family of Sanskrit \emph{vighaṭate}, \textquotedblleft{}flies apart, is broken\textquotedblright{}, where CDIAL lists Marathi \emph{bighaḍṇẽ}, \textquotedblleft{}to be spoilt\textquotedblright{} (CDIAL 11673).
\end{cousinweb}

\begin{grammarlens}[title={What you've built}]
One more word for getting around.
\end{grammarlens}

\subsection*{Your turn}
\begin{itemize}
\raggedright
  \item \emph{Say it:} \emph{bighāḍ}
  \item \emph{Say it:} \emph{bighāḍ}, once more
  \item \emph{Say it:} \emph{iñjināt bighāḍ jhālyāmuḷe gāḍī thāmblī.} — The train stopped because of a fault in the engine.
  \item \emph{Recall:} say the Marathi for a share, then the Marathi for to be left over, then say \emph{bighāḍ} again
\end{itemize}

\begin{tcolorbox}[breakable,colback=teal!4,colframe=teal!35!black,title={Before you move on}]
What does \mr{बिघाड} mean? (\textquotedblleft{}A breakdown; a fault\textquotedblright{}.) Say it once more.
\end{tcolorbox}
\section[\texorpdfstring{\mr{निकामी} (nikāmī)}{निकामी (nikāmī)}]{\mr{निकामी} (nikāmī) — out of action; useless}

\label{lesson:MR-C431-nikami}



\begin{quote}
Before the new one: say the Marathi for to be left over, then the Marathi for a breakdown.
\end{quote}

\subsection*{You'll want to know: \mr{निकामी}}
\textbf{\mr{निकामी}} — \emph{nikāmī} — \textquotedblleft{}out of action; useless\textquotedblright{}.

\begin{quote}
\mr{स्टेशनवरचा} \mr{जुना} \mr{पंखा} \mr{आता} \mr{निकामी} \mr{झाला} \mr{आहे}.
\end{quote}

\begin{quote}
\emph{sṭeśanvarcā junā paṅkhā ātā nikāmī jhālā āhe.}
\end{quote}

\begin{quote}
The old fan at the station no longer works at all.
\end{quote}

It does not change: \mr{निकामी} \mr{पंखा}, \mr{निकामी} \mr{गाडी}; \mr{निकामी} \mr{होणे}, to stop working for good.

\begin{cousinweb}
\textbf{\mr{निकामी}} (\emph{nikāmī}) is inherited: CDIAL lists Marathi \emph{nikāmī}, \textquotedblleft{}idle, useless\textquotedblright{}, under Sanskrit \emph{niṣkarman}, \textquotedblleft{}inactive\textquotedblright{}, that is, without work (CDIAL 7475).
\end{cousinweb}

\begin{grammarlens}[title={What you've built}]
One more word for getting around.
\end{grammarlens}

\subsection*{Your turn}
\begin{itemize}
\raggedright
  \item \emph{Say it:} \emph{nikāmī}
  \item \emph{Say it:} \emph{nikāmī}, once more
  \item \emph{Say it:} \emph{sṭeśanvarcā junā paṅkhā ātā nikāmī jhālā āhe.} — The old fan at the station no longer works at all.
  \item \emph{Recall:} say the Marathi for to be left over, then the Marathi for a breakdown, then say \emph{nikāmī} again
\end{itemize}

\begin{tcolorbox}[breakable,colback=teal!4,colframe=teal!35!black,title={Before you move on}]
What does \mr{निकामी} mean? (\textquotedblleft{}Out of action; useless\textquotedblright{}.) Say it once more.
\end{tcolorbox}
\section[\texorpdfstring{\mr{तुटणे} (tuṭṇe)}{तुटणे (tuṭṇe)}]{\mr{तुटणे} (tuṭṇe) — to snap; to break}

\label{lesson:MR-C431-tutne}



\begin{quote}
Before the new one: say the Marathi for a breakdown, then the Marathi for out of action.
\end{quote}

\subsection*{You'll want to know: \mr{तुटणे}}
\textbf{\mr{तुटणे}} — \emph{tuṭṇe} — \textquotedblleft{}to snap; to break\textquotedblright{}.

\begin{quote}
\mr{वादळात} \mr{झाडाची} \mr{फांदी} \mr{तुटली}.
\end{quote}

\begin{quote}
\emph{vādaḷāt jhāḍācī phāndī tuṭlī.}
\end{quote}

\begin{quote}
The branch of the tree snapped in the storm.
\end{quote}

Past (he …): \mr{तो} \mr{तुटला}, \emph{to tuṭlā}.

It happens by itself: \mr{फांदी} \mr{तुटली}; to break something is \mr{तोडणे}.

\begin{cousinweb}
\textbf{\mr{तुटणे}} (\emph{tuṭṇe}) is inherited: CDIAL lists Marathi \emph{tuṭṇẽ}, \textquotedblleft{}is broken\textquotedblright{}, under Sanskrit \emph{truṭyati}, \textquotedblleft{}is broken, falls asunder\textquotedblright{} (CDIAL 6065).
\end{cousinweb}

\begin{grammarlens}[title={What you've built}]
One more word for getting around.
\end{grammarlens}

\subsection*{Your turn}
\begin{itemize}
\raggedright
  \item \emph{Say it:} \emph{tuṭṇe}
  \item \emph{Say it:} \emph{tuṭṇe}, once more
  \item \emph{Say it:} \emph{vādaḷāt jhāḍācī phāndī tuṭlī.} — The branch of the tree snapped in the storm.
  \item \emph{Recall:} say the Marathi for a breakdown, then the Marathi for out of action, then say \emph{tuṭṇe} again
\end{itemize}

\begin{tcolorbox}[breakable,colback=teal!4,colframe=teal!35!black,title={Before you move on}]
What does \mr{तुटणे} mean? (\textquotedblleft{}To snap; to break\textquotedblright{}.) Say it once more.
\end{tcolorbox}
\section[\texorpdfstring{\mr{फुटणे} (phuṭṇe)}{फुटणे (phuṭṇe)}]{\mr{फुटणे} (phuṭṇe) — to burst}

\label{lesson:MR-C431-phutne}



\begin{quote}
Before the new one: say the Marathi for out of action, then the Marathi for to snap.
\end{quote}

\subsection*{You'll want to know: \mr{फुटणे}}
\textbf{\mr{फुटणे}} — \emph{phuṭṇe} — \textquotedblleft{}to burst\textquotedblright{}.

\begin{quote}
\mr{वाटेत} \mr{गाडीचा} \mr{टायर} \mr{फुटला}.
\end{quote}

\begin{quote}
\emph{vāṭet gāḍīcā ṭāyar phuṭlā.}
\end{quote}

\begin{quote}
On the way a tyre of the car burst.
\end{quote}

Past (he …): \mr{तो} \mr{फुटला}, \emph{to phuṭlā}.

Also a crack in glass or a pot: \mr{काच} \mr{फुटली}, the glass broke.

\begin{cousinweb}
\textbf{\mr{फुटणे}} (\emph{phuṭṇe}) is inherited: CDIAL lists Marathi \emph{phuṭṇẽ}, \textquotedblleft{}to split, burst\textquotedblright{}, under \emph{*sphuṭyati}, \textquotedblleft{}is burst open\textquotedblright{} (CDIAL 13845).
\end{cousinweb}

\begin{grammarlens}[title={What you've built}]
One more word for getting around.
\end{grammarlens}

\subsection*{Your turn}
\begin{itemize}
\raggedright
  \item \emph{Say it:} \emph{phuṭṇe}
  \item \emph{Say it:} \emph{phuṭṇe}, once more
  \item \emph{Say it:} \emph{vāṭet gāḍīcā ṭāyar phuṭlā.} — On the way a tyre of the car burst.
  \item \emph{Recall:} say the Marathi for out of action, then the Marathi for to snap, then say \emph{phuṭṇe} again
\end{itemize}

\begin{tcolorbox}[breakable,colback=teal!4,colframe=teal!35!black,title={Before you move on}]
What does \mr{फुटणे} mean? (\textquotedblleft{}To burst\textquotedblright{}.) Say it once more.
\end{tcolorbox}
\section[\texorpdfstring{\mr{फाटणे} (phāṭṇe)}{फाटणे (phāṭṇe)}]{\mr{फाटणे} (phāṭṇe) — to rip; to tear}

\label{lesson:MR-C431-phatne}



\begin{quote}
Before the new one: say the Marathi for to snap, then the Marathi for to burst.
\end{quote}

\subsection*{You'll want to know: \mr{फाटणे}}
\textbf{\mr{फाटणे}} — \emph{phāṭṇe} — \textquotedblleft{}to rip; to tear\textquotedblright{}.

\begin{quote}
\mr{गर्दीत} \mr{माझी} \mr{पिशवी} \mr{फाटली}.
\end{quote}

\begin{quote}
\emph{gardīt mājhī piśvī phāṭlī.}
\end{quote}

\begin{quote}
My bag ripped in the crowd.
\end{quote}

Past (he …): \mr{तो} \mr{फाटला}, \emph{to phāṭlā}.

Cloth and paper \mr{फाटतात}; to tear something is \mr{फाडणे}.

\begin{cousinweb}
\textbf{\mr{फाटणे}} (\emph{phāṭṇe}) is inherited: CDIAL lists Marathi \emph{phāṭṇẽ}, \textquotedblleft{}to be split\textquotedblright{}, under Sanskrit \emph{sphāṭayati}, \textquotedblleft{}splits\textquotedblright{} (CDIAL 13825).
\end{cousinweb}

\begin{grammarlens}[title={What you've built}]
One more word for getting around.
\end{grammarlens}

\subsection*{Your turn}
\begin{itemize}
\raggedright
  \item \emph{Say it:} \emph{phāṭṇe}
  \item \emph{Say it:} \emph{phāṭṇe}, once more
  \item \emph{Say it:} \emph{gardīt mājhī piśvī phāṭlī.} — My bag ripped in the crowd.
  \item \emph{Recall:} say the Marathi for to snap, then the Marathi for to burst, then say \emph{phāṭṇe} again
\end{itemize}

\begin{tcolorbox}[breakable,colback=teal!4,colframe=teal!35!black,title={Before you move on}]
What does \mr{फाटणे} mean? (\textquotedblleft{}To rip; to tear\textquotedblright{}.) Say it once more.
\end{tcolorbox}
